\documentclass[11pt,a4paper]{article}

\usepackage[T1]{fontenc}
\usepackage[utf8]{inputenc}
\usepackage[ngerman]{babel}
\usepackage[a4paper,margin=2.4cm]{geometry}
\usepackage{amsmath,amssymb}
\usepackage{booktabs}
\usepackage{tabularx}
\usepackage{array}
\usepackage{xcolor}
\usepackage{enumitem}
\usepackage{fancyhdr}
\usepackage[hidelinks]{hyperref}

% Spaltentypen
\newcolumntype{L}{>{\raggedright\arraybackslash}X}
\newcolumntype{C}{>{\centering\arraybackslash}X}
\newcolumntype{r}{>{\raggedleft\arraybackslash}p{1.9cm}}

% Kopf-/Fusszeile
\pagestyle{fancy}
\fancyhf{}
\fancyhead[L]{\small Q-Learning -- Analyseprotokoll Lauf 2}
\fancyhead[R]{\small Symmetrienutzung}
\fancyfoot[C]{\small\thepage}

% Statusmarker
\newcommand{\ok}{\textcolor{green!55!black}{\textbf{ok}}}
\newcommand{\bad}{\textcolor{red!75!black}{\textbf{Fehler}}}
\newcommand{\offen}{\textcolor{orange!85!black}{\textbf{offen}}}

\title{\textbf{Analyseprotokoll Lauf 2}\\[0.35em]
\large Tabellarisches Q-Learning zur robotergestützten Sortierung\\
\large mit Ausnutzung der Puffersymmetrie}
\author{Niko \\ \small Diskrete Simulation und Reinforcement Learning}
\date{\today}

\begin{document}

\maketitle
\thispagestyle{fancy}

\begin{abstract}
\noindent
Dieses Protokoll dokumentiert den zweiten Trainingslauf des tabellarischen
Q-Learning-Agenten für die Sortierung eingehender Teile auf zwei Puffer mit
Return-Loop. Gegenüber Lauf 1 wurde ausschließlich die Maßnahme M4
(Ausnutzung der Puffersymmetrie) umgesetzt. Die Symmetrie ist bit-exakt
korrekt implementiert; die Zahl der unterschiedlichen Fehlerlogiken sinkt von
acht auf fünf, und alle vier explizit dokumentierten Fehlerzeilen aus Lauf 1
sind behoben. Das Explorationsdefizit (M1--M3) wurde nicht adressiert und
äußert sich in Lauf 2 an neuen, teils gravierenderen Stellen. Als neuer
Hauptbefund wird die weiterhin ungenutzte \emph{Typ}-Symmetrie
identifiziert.
\end{abstract}

\tableofcontents
\bigskip

%======================================================================
\section{Zielsetzung und Änderungen gegenüber Lauf 1}
%======================================================================

Ziel von Lauf 2 war die Überprüfung der Wirksamkeit von Maßnahme M4 aus dem
Protokoll zu Lauf 1: die Ausnutzung der strukturellen Symmetrie der beiden
Puffer. Die Spiegelung ist definiert als der gleichzeitige Tausch von
\[
  (b1typ,\,b1cat) \;\longleftrightarrow\; (b2typ,\,b2cat)
  \qquad\text{und}\qquad
  \texttt{Action\_0} \;\longleftrightarrow\; \texttt{Action\_1}.
\]

\begin{table}[htbp]
\centering
\caption{Umsetzungsstand der in Lauf 1 empfohlenen Maßnahmen}
\label{tab:massnahmen-status}
\begin{tabularx}{\textwidth}{@{}clLc@{}}
\toprule
\textbf{Nr.} & \textbf{Maßnahme} & \textbf{Zweck} & \textbf{Status} \\
\midrule
M1 & Optimistische $Q$-Initialisierung & Erzwingt Erstbesuch aller Aktionen & \offen \\
M2 & $\varepsilon$-Minimum & Verhindert vollständigen Explorationsstopp & \offen \\
M3 & UCB / Besuchszähler & Gezielte Exploration seltener Zustände & \offen \\
M4 & Puffersymmetrie & Halbierung der zu lernenden Klassen & \ok \\
M5 & Kompaktere Zustandskodierung & Eliminiert unmögliche Zustände & \offen \\
M6 & Nächster Ankunftstyp im Zustand & Ermöglicht Lookahead & \offen \\
\bottomrule
\end{tabularx}
\end{table}

\noindent
Reward-Design, Diskontfaktor und Zustandsdimension blieben unverändert. Die
Q-Tabelle umfasst weiterhin $432$ Zustände $\times$ $3$ Aktionen
(Puffer~1, Puffer~2, Return-Loop).

%======================================================================
\section{Verifikation der Symmetrieimplementierung}
%======================================================================

\subsection{Spiegelpaare}

Für Spiegelpaare muss der $Q$-Wert der jeweils gewählten Aktion exakt
übereinstimmen. Tabelle~\ref{tab:spiegelpaare} zeigt eine repräsentative
Stichprobe.

\begin{table}[htbp]
\centering
\caption{Stichprobe der Spiegelpaare, Zustand als $(inc,\,b1typ,\,b2typ,\,b1cat,\,b2cat)$}
\label{tab:spiegelpaare}
\small
\begin{tabular}{@{}llr lr c@{}}
\toprule
\textbf{Paar} & \textbf{Zustand} & \textbf{$Q$ gewählt} &
\textbf{Gespiegelt} & \textbf{$Q$ gewählt} & \textbf{Status} \\
\midrule
9 / 36    & (1,0,1,1,1) & 417,64688454714667 & (1,1,0,1,1) & 417,64688454714667 & \ok \\
54 / 81   & (1,1,2,1,1) & 412,9399113166976  & (1,2,1,1,1) & 412,9399113166976  & \ok \\
99 / 126  & (1,2,3,1,1) & 402,826410225901   & (1,3,2,1,1) & 402,826410225901   & \ok \\
359 / 413 & (3,1,3,3,3) & 510,6411597942206  & (3,3,1,3,3) & 510,6411597942206  & \ok \\
\bottomrule
\end{tabular}
\end{table}

\subsection{Selbstsymmetrische Zustände}

Zustände mit $b1typ = b2typ$ \emph{und} $b1cat = b2cat$ sind zu sich selbst
symmetrisch; für sie muss
$Q(s,\texttt{Action\_0}) = Q(s,\texttt{Action\_1})$ gelten. Dies ist in
Lauf 2 durchgängig erfüllt für die Zeilen
\[
  0,\;45,\;90,\;135,\;242,\;279,\;283,\;333,\;378,\;386,\;423 .
\]
In Lauf 1 wählten die Leerzustände $0$, $144$ und $288$ noch widersprüchlich
Puffer~2, Puffer~1 und Puffer~1. Diese Inkonsistenz ist behoben.

\begin{center}
\fbox{\parbox{0.93\textwidth}{\textbf{Bewertung:} Maßnahme M4 ist
vollständig und fehlerfrei umgesetzt. Effektiv mussten nur
\textbf{111 Symmetrieklassen} statt 219 betriebsrelevanter Zustände gelernt
werden; die Sample-Effizienz verdoppelt sich damit rechnerisch.}}
\end{center}

%======================================================================
\section{Abdeckung des Zustandsraums}
%======================================================================

\begin{table}[htbp]
\centering
\caption{Aufteilung der 432 Zeilen der Q-Tabelle}
\label{tab:abdeckung}
\begin{tabularx}{\textwidth}{@{}Lrrl@{}}
\toprule
\textbf{Kategorie} & \textbf{Zustände} & \textbf{Anteil} & \textbf{Status} \\
\midrule
Strukturell unmöglich ($typ=0 \wedge cat>1$) & 132 & 30,6\,\% & alle exakt $0{,}0$ -- korrekt \\
Betriebsrelevant                             & 219 & 50,7\,\% & 100\,\% trainiert, $Q \approx 387\ldots511$ \\
Beide Puffer gleicher Typ                    &  81 & 18,8\,\% & untrainiert, alle $Q < 81$ \\
\quad davon nie besucht                      &  28 &  6,5\,\% & exakt $0{,}0$ \\
\midrule
\textbf{Summe}                               & \textbf{432} & \textbf{100\,\%} & \\
\bottomrule
\end{tabularx}
\end{table}

\noindent
Die 132 unmöglichen Zeilen existieren weiterhin, da Maßnahme M5 nicht
umgesetzt wurde. Sie belegen unverändert $30{,}6\,\%$ des Speichers, ohne
Information zu tragen.

%======================================================================
\section{Qualität der gelernten Policy}
%======================================================================

\begin{table}[htbp]
\centering
\caption{Trefferquoten nach Entscheidungssituation, Lauf 2 gegen Lauf 1}
\label{tab:policy}
\small
\begin{tabularx}{\textwidth}{@{}LlrrR{}}
\toprule
\textbf{Situation} & \textbf{Korrekte Aktion} & \textbf{Treffer} &
\textbf{Quote} & \textbf{Lauf 1} \\
\midrule
Beide Puffer leer & symmetrisch & 3 / 3 & --- & inkonsistent \\
Ein Puffer leer, anderer = $inc$ (Match) & Match-Puffer & 18 / 18 & \textbf{100,0\,\%} & 95--97\,\% \\
Ein Puffer leer, anderer Fremdtyp & leeren Puffer belegen & 32 / 36 & 88,9\,\% & 91,7\,\% \\
Beide belegt, ein Match & Match-Puffer & 102 / 108 & 94,4\,\% & $\approx$ 96\,\% \\
Beide belegt, kein Match (strikt) & Return & 34 / 54 & 63,0\,\% & 68,5\,\% \\
\quad dito, Opfern als valide gewertet & --- & 54 / 54 & \textbf{100,0\,\%} & 9 von 11 \\
\midrule
\textbf{Gesamt, strikte Bewertung} & & \textbf{189 / 219} & \textbf{86,3\,\%} & 88,4\,\% \\
\textbf{Gesamt, mit Opfer-Toleranz} & & \textbf{209 / 219} & \textbf{95,4\,\%} & 92,6\,\% \\
\bottomrule
\end{tabularx}
\end{table}

\subsection*{Interpretation}

\begin{itemize}[leftmargin=1.4em]
  \item Der \textbf{Match-Vorrang bei leerem Gegenpuffer} ist mit
        $100\,\%$ vollständig gelernt.
  \item Das \textbf{Opfer-Kriterium} ist nun strikt konsistent: In allen
        20 Opferentscheidungen wird ausnahmslos der \emph{weniger} gefüllte
        Puffer überschrieben, nie der vollere. In Lauf 1 galt dies nur in
        9 von 11 Fällen.
  \item Die \textbf{strikte Gesamtquote sinkt} leicht von $88{,}4\,\%$ auf
        $86{,}3\,\%$. Dies ist ein Artefakt der Symmetrie: Fehler treten
        zwangsläufig paarweise auf und werden auf Zustandsebene doppelt
        gezählt, obwohl es sich um dieselbe Entscheidungslogik handelt.
        Die tolerante Quote steigt gleichzeitig von $92{,}6\,\%$ auf
        $95{,}4\,\%$.
\end{itemize}

%======================================================================
\section{Verbleibende Fehlerklassen}
%======================================================================

\begin{table}[htbp]
\centering
\caption{Die fünf verbleibenden Fehlerlogiken (je zwei gespiegelte Zeilen)}
\label{tab:fehler}
\small
\begin{tabular}{@{}llrlr@{}}
\toprule
\textbf{Klasse} & \textbf{Zustand (Spiegel)} & \textbf{Gewählt} &
\textbf{$Q$ gewählt} & \textbf{Korrekt / $Q$ korrekt} \\
\midrule
K1 & 20 / 78 \; (1,0,2,1,\textbf{3})            & Return & 445,01 & P1 \; / \; \textbf{0,092} \\
K2 & 298 / 327 \; (3,0,1,1,2)                   & Return & 421,93 & P1 \; / \; \textbf{$-0{,}035$} \\
K3 & 200 / 231 \; (2,1,2,1,\textbf{3})          & P1     & 420,59 & P2 (Match!) \; / \; \textbf{2,198} \\
K4 & 206 / 233 \; (2,1,2,\textbf{3},\textbf{3}) & Return & 448,08 & P2 (Match!) \; / \; \textbf{2,042} \\
K5 & 251 / 278 \; (2,2,\textbf{3},\textbf{3},\textbf{3}) & P2 & 423,17 & P1 (Match!) \; / \; \textbf{3,671} \\
\bottomrule
\end{tabular}
\end{table}

\subsection{Diagnose}

In allen fünf Klassen liegt der $Q$-Wert der korrekten Aktion unter $4$,
entspricht also null bis zwei Besuchen. Der \texttt{argmax} entscheidet
damit gegen praktisch uninitialisierte Werte. Vier der fünf Klassen
enthalten die seltene Füllstandskategorie $cat=3$. Die Ursache ist somit
\textbf{unverändert ein Explorations- und $\varepsilon$-Decay-Problem} und
\emph{kein} Reward- oder Regelproblem.

\subsection{Verschärfung gegenüber Lauf 1}

Die Fehler haben ihre Qualität verändert:

\begin{itemize}[leftmargin=1.4em]
  \item In \textbf{Lauf 1} lagen alle acht Fehler in der ökonomisch
        günstigeren Kategorie „leerer Puffer versus Return``.
  \item In \textbf{Lauf 2} betreffen K3--K5 nun \emph{Match-Zustände mit}
        $cat=3$. In K5 überschreibt der Agent einen zu etwa $90\,\%$
        gefüllten Fremdpuffer, statt seinen eigenen nahezu vollständigen
        Stapel zu vollenden. Das kostet zwei Auslagerungen statt keiner.
\end{itemize}

\noindent
Die alten Fehlerzeilen 29, 61, 173 und 299 aus Lauf 1 sind sämtlich
korrigiert. Die Zahl der unterschiedlichen Fehlerlogiken sinkt von
\textbf{8 auf 5}.

%======================================================================
\section{Neuer Hauptbefund: ungenutzte Typ-Symmetrie}
%======================================================================

Die \emph{Puffer} sind nach M4 symmetrisch, die \emph{Typ-Labels} jedoch
nicht. Funktional identische Situationen werden weiterhin sechsfach
getrennt gelernt.

\begin{table}[htbp]
\centering
\caption{Sechs funktional identische Zustände: Match bei $cat=1$, Fremdpuffer $cat=1$}
\label{tab:typredundanz}
\begin{tabularx}{0.82\textwidth}{@{}rLr@{}}
\toprule
\textbf{Zustand} & \textbf{Bedeutung} & \textbf{$Q$ (Match-Aktion)} \\
\midrule
54  & Match $cat=1$, Fremd $cat=1$ & 412,94 \\
63  & dito & 430,95 \\
225 & dito & 402,16 \\
243 & dito & 404,65 \\
405 & dito & 434,26 \\
414 & dito & 407,67 \\
\midrule
\multicolumn{2}{@{}l}{\textbf{Spannweite}} & \textbf{402--434 ($\approx 8\,\%$)} \\
\bottomrule
\end{tabularx}
\end{table}

\noindent
Bei $cat=3$ wird die Streuung kritisch: Die Werte lauten
$494{,}85$, $496{,}70$, $510{,}64$ und $440{,}79$ -- und zwei der sechs
Fälle (K4, K5) sind vollständig falsch gelernt. Inhaltlich zeigt sich die
Redundanz ebenso: Zustand 105 $(3,2,3,3,1)$ opfert korrekt den
$cat=1$-Puffer, der typgleiche Fall 348 $(3,1,2,3,1)$ wählt hingegen Return.

\subsection{Lösungsvorschlag: relative Kodierung}

Eine relative Kodierung je Puffer über
$\{\text{leer},\,\text{match},\,\text{mismatch}\} \times cat$
reduziert den Zustandsraum auf $7 \times 7 = 49$ und mit Puffersymmetrie auf
\textbf{28 Klassen}. Sowohl $inc$ als auch die Typ-Labels entfallen
vollständig, und die Daten aller sechs Blöcke werden gepoolt. Bei etwa
sechsfacher effektiver Datenmenge pro Klasse ist die Beseitigung von
K3--K5 mit hoher Wahrscheinlichkeit zu erwarten.

%======================================================================
\section{Struktur der Wertefunktion}
%======================================================================

\begin{table}[htbp]
\centering
\caption{$Q$-Werte der gewählten Aktion, Zustände 351--359
($inc=3$, $b1typ=1$, $b2typ=3$, Match auf Puffer 2)}
\label{tab:wertefunktion}
\begin{tabular}{@{}c rrr@{}}
\toprule
$b1cat \,\backslash\, b2cat$ & \textbf{1} & \textbf{2} & \textbf{3} \\
\midrule
\textbf{1} & 434,26 & 462,30 & 492,39 \\
\textbf{2} & 452,33 & 472,78 & 496,32 \\
\textbf{3} & 469,14 & 483,14 & \textbf{510,64} \\
\bottomrule
\end{tabular}
\end{table}

\noindent
Die Werte sind streng monoton steigend in \emph{beiden} Füllständen und
damit sauberer als in Lauf 1, wo Dubletten wie $669{,}2$ auftraten. Die
Füllstandsmerkmale werden also nachweislich informativ genutzt. In anderen
Blöcken (etwa 54--62) ist die Monotonie allerdings nur bezüglich des
\emph{Zielpuffers} stabil, bezüglich des Gegenpuffers verrauscht.

\medskip\noindent
\textbf{Unverändert problematisch:} Die nicht gewählten Aktionen liegen bei
$-0{,}1 \ldots -2$ und tragen keine Information. Es existiert weiterhin
\textbf{keine verwertbare zweitbeste Aktion}, was Fallback-Strategien und
Policy-Robustheitsanalysen unmöglich macht.

%======================================================================
\section{Auffälligkeit: Absenkung des \texorpdfstring{$Q$}{Q}-Niveaus}
%======================================================================

\begin{table}[htbp]
\centering
\caption{Vergleich der absoluten $Q$-Niveaus}
\label{tab:qniveau}
\begin{tabularx}{0.85\textwidth}{@{}Lrrr@{}}
\toprule
 & \textbf{Lauf 1} & \textbf{Lauf 2} & \textbf{Verhältnis} \\
\midrule
Maximum              & 700,2 & 510,64 & 72,9\,\% \\
Minimum (trainiert)  & 509,5 & 386,87 & 75,9\,\% \\
\bottomrule
\end{tabularx}
\end{table}

\noindent
Bei identischem Reward-Design und identischer Episodenzahl müsste die
Symmetrienutzung die Werte \emph{schneller} konvergieren lassen, nicht auf
etwa $73\,\%$ absenken.

\begin{center}
\fbox{\parbox{0.93\textwidth}{\textbf{Offene Frage:} Wurden Reward,
$\gamma$, Episodenzahl oder Episodenlänge mitverändert? Falls nicht, deutet
das Niveau auf eine \emph{unterkonvergierte} Wertefunktion hin. In diesem
Fall sind die Quoten in Tabelle~\ref{tab:policy} als untere Schranke zu
lesen.}}
\end{center}

%======================================================================
\section{Maßnahmen für Lauf 3}
%======================================================================

\begin{table}[htbp]
\centering
\caption{Priorisierte Maßnahmenliste}
\label{tab:massnahmen3}
\begin{tabularx}{\textwidth}{@{}clLl@{}}
\toprule
\textbf{Nr.} & \textbf{Maßnahme} & \textbf{Bezug} & \textbf{Priorität} \\
\midrule
M1 & Optimistische Initialisierung $Q_0 \approx 520$ & K1--K5 & \textbf{hoch} \\
M2 & $\varepsilon_{\min} = 0{,}05$ statt vollständigem Decay & K1--K5 & \textbf{hoch} \\
M3 & Besuchszählbasierte Exploration (UCB) & K3--K5 & \textbf{hoch} \\
M8 & Relative Typkodierung $\{$leer, match, mismatch$\}$ & Abschnitt 6 & \textbf{hoch} \\
M5 & $cat=0$ für leere Puffer & 132 Leerzeilen & mittel \\
M7 & Besuchszähler mitloggen & Auswertbarkeit & mittel \\
M6 & Nächsten Ankunftstyp in den Zustand & Return vs. Opfer & niedrig \\
\bottomrule
\end{tabularx}
\end{table}

\noindent
Maßnahme M4 ist erledigt und muss nicht erneut angefasst werden. Ohne
M1--M3 wird jeder weitere Lauf dieselben drei bis acht unerforschten
$cat=3$-Zustände falsch entscheiden -- lediglich an wechselnden Stellen.

%======================================================================
\section{Fazit}
%======================================================================

\begin{enumerate}[leftmargin=1.6em]
  \item Die Puffersymmetrie ist \textbf{bit-exakt korrekt} implementiert;
        alle Spiegelpaare und selbstsymmetrischen Zustände sind konsistent.
  \item Die Zahl der unterschiedlichen Fehlerlogiken sinkt von
        \textbf{8 auf 5}; alle vier dokumentierten Fehlerzeilen aus Lauf 1
        sind behoben.
  \item Match-Vorrang bei leerem Gegenpuffer ($100\,\%$) und
        Opfer-Kriterium ($100\,\%$ konsistent) sind vollständig gelernt.
  \item Das \textbf{Explorationsdefizit besteht unverändert} und trifft in
        Lauf 2 sogar ökonomisch teurere Match-Zustände mit $cat=3$.
  \item Als neuer Hauptbefund ist die \textbf{Typ-Symmetrie} zu nennen:
        sechsfach redundantes Lernen identischer Situationen mit bis zu
        $8\,\%$ Wertestreuung. Eine relative Kodierung würde den
        Zustandsraum auf 28 Klassen reduzieren.
  \item Das abgesenkte $Q$-Niveau ($\approx 73\,\%$ von Lauf 1) ist vor der
        Interpretation der Quoten zu klären.
\end{enumerate}

\end{document}